\section{Consistency checks}\label{S:consistency}

These comparisons place the second-moment theorem alongside known bounds and
Chen's conjecture. They are not inputs to the proof of
Theorem~\ref{T:intro-second-moment}.

\paragraph{Time-independent noise.}
As $\alpha_0\downarrow0$, the exponent $\kappa$ tends to
$(4-\alpha)/(3-\alpha)$, the exponent of the time-independent Stratonovich
model found by Chen and Hu~\cite{chen.hu:25:hyperbolic}*{Theorem~1.2}. This is
a formal boundary limit: the large-time limit and the limit
$\alpha_0\downarrow0$ are not interchanged (Remark~\ref{R:G1-scope}).

\paragraph{The first moment.}
If $\liminf_{t\to\infty}t^{-\kappa}\log\E u(t,x)\ge\ell$, then Jensen's
inequality $\E u^2(t,x)\ge(\E u(t,x))^2$ gives the corresponding second-moment
lower limit at least $2\ell$ (Proposition~\ref{P:Jensen}). Hence
$C_2(\Gamma,\alpha_0)\ge2\ell$ in Theorem~\ref{T:intro-second-moment}. Taking
$\ell=L_*$ of Theorem~\ref{T:first} is conditional on the intended first-moment
lower bound with that constant; its printed normalization is not adopted
unconditionally.

\paragraph{The Skorohod model.}
For Riesz covariances with $a=\alpha<d$ and $\alpha<2$, set
$H=1-\alpha_0/2>1/2$. The Skorohod upper bound of Balan and
Conus~\cite{balan.conus:16:intermittency}*{Theorem~2.1(b)} has the time
exponent $(2H+2-a)/(3-a)=\kappa$. The comparison concerns time exponents only:
the solution concepts and covariance normalizations differ, and no constants
are compared. The Stratonovich moment expansion also includes self-pairings,
that is, pairs of noise factors from the same ordered integral, which are
absent from the Skorohod expansion (Remark~\ref{R:G1-scope}).

\paragraph{Chen's conjecture.}
For $p=1,2,\dots$, Chen conjectures that $t^{-\kappa}\log\E|u(t,x)|^p$ tends to
$p^{(4-\alpha)/(3-\alpha)}L^*$ (Conjecture~\ref{Conj:allp}). For $p=2$ the
conjectured value is $2^{(4-\alpha)/(3-\alpha)}L^*$. When $L^*>0$, this exceeds
$2L^*$, the lower bound that Jensen's inequality would give if the first-moment
rate were $L^*$; at zero covariance both values vanish.
Theorem~\ref{T:intro-second-moment} does not decide this conjecture: its
constant $C_2(\Gamma,\alpha_0)$ is given by its own variational formula and is
not identified with the conjectured constant or with the first-moment limit.
The conjectured value is a point of comparison, not a target of the proof.
